\chapter{EPI Ledger --- Epistemology Tags}
\label{chap:EPI_Ledger}

\section{Purpose}

The EPI ledger classifies Engrams according to epistemic status: how knowledge
is justified, how claims are made, and what evidential structures exist.

\section{Scope}

Includes:
\begin{itemize}
    \item axioms,
    \item inference rules,
    \item empirical evidence,
    \item logical forms,
    \item certainty/credence markers.
\end{itemize}

\section{Schema}

\[
T_{\text{EPI}} = \langle key,\, gloss,\, bindings,\, constraints \rangle
\]

Examples:
\begin{itemize}
    \item \texttt{epi:axiom}
    \item \texttt{epi:empirical}
    \item \texttt{epi:inference}
    \item \texttt{epi:evidence}
    \item \texttt{epi:confidence}
\end{itemize}

\section{Class Binding Notes}

\begin{itemize}
    \item S-class for structural reasoning or logic.
    \item C-class for validated canonical claims.
    \item O-class for narrative epistemic uncertainty.
\end{itemize}

\section{Constraints}

\begin{itemize}
    \item Epistemic statuses must match content of Engram.
    \item Claims requiring high confidence must not contradict ZED.
    \item Axioms must be P-class reviewed before kernel use.
\end{itemize}
